\documentclass[12pt]{article}
\usepackage[margin=0.3in,top=0.25in,bottom=0.4in]{geometry}
\makeatletter
\renewcommand{\@maketitle}{\newpage\null\vskip -3.5em\begin{center}{\LARGE \@title \par}\end{center}\par\vskip 0.2em}
\makeatother
\usepackage{amsmath,amssymb}
\usepackage{booktabs}
\usepackage{hyperref}
\setlength{\parskip}{0.6em}
\setlength{\parindent}{0pt}
\newcommand{\xpo}{x_{\mathbb{P}}}

\title{How $R_{pA}$ scales with $A$: exclusive vs.\ diffractive (TMD)}
\author{}
\date{}

\begin{document}
\maketitle

\section*{In short}

We define
\[
  R_{pA}=\frac{\sigma_A}{A\,\sigma_p}\,,
\]
so a process whose cross section grows like $A$ has $R_{pA}=1$. The two contributions behave very
differently:
\begin{itemize}
\item \textbf{Exclusive:} at large $p_\perp$ it grows like $A^{4/3}$, so $R_{pA}\propto A^{1/3}$. For Pb
  this gives $R_{pA}\approx2$--$2.5$, as in our figures.
\item \textbf{Diffractive (TMD):} it grows like $A$, so $R_{pA}\to$ constant $\approx1$ (about $0.9$ in our
  figures), even at large $p_\perp$.
\end{itemize}
The reason is that the exclusive process at large $p_\perp$ probes small, dilute dipoles, while the
diffractive TMD is always dominated by gluons with $K_\perp\sim Q_s$, i.e.\ by the saturation region.

\section*{1. Exclusive: $R_{pA}\propto A^{1/3}$}

At large $p_\perp$ the radial integrals $\mathcal{Q}_{0,1}$ of Eqs.~(2)--(3) are dominated by small dipoles,
$r_\perp\sim1/p_\perp$. For small dipoles the nucleus is dilute, and the nuclear dipole amplitude is
the proton one multiplied by the number of nucleons met along the path (optical Glauber):
\[
  \mathcal{N}_A(r_\perp,\mathbf{b})\approx A\,T_A(\mathbf{b})\,\frac{\sigma_0}{2}\,\mathcal{N}_p(r_\perp)\,,
\]
where $T_A$ is the nuclear thickness function ($\int\mathrm{d}^2\mathbf{b}\,T_A=1$) and $\sigma_0/2$ is the
transverse area of the proton.

\textbf{Where this comes from.} The dipole scatters independently off each nucleon it meets:
\begin{itemize}
\item \emph{One proton.} The proton is a disk of area $\sigma_0/2$, and a dipole that hits it scatters
with amplitude $\mathcal{N}_p(r_\perp)$. The dipole--proton cross section is
$\sigma^p_{\rm dip}=2\,(\sigma_0/2)\,\mathcal{N}_p$.
\item \emph{Nucleons along the path.} $A\,T_A(\mathbf{b})$ is the number of nucleons per unit area at
impact parameter $\mathbf{b}$. Times the proton area $\sigma_0/2$, it is the number of nucleons that
the dipole meets.
\item \emph{No scattering at all.} The amplitude to cross the nucleus without scattering is the product
over the nucleons, which exponentiates:
\[
  1-\mathcal{N}_A(r_\perp,\mathbf{b})=\Big[1-T_A(\mathbf{b})\,\frac{\sigma_0}{2}\,\mathcal{N}_p(r_\perp)\Big]^A
  \simeq\exp\Big[-A\,T_A(\mathbf{b})\,\frac{\sigma_0}{2}\,\mathcal{N}_p(r_\perp)\Big].
\]
\item \emph{Small dipole.} When the exponent is small, expanding gives the formula above: the number of
nucleons met times the amplitude on one of them. When it is not small, $\mathcal{N}_A$ is smaller than
this sum (saturation).
\end{itemize}
\emph{Check with the tables.} At the center of Pb, $A\,T_A(0)\approx2.2$~fm$^{-2}$ and
$\sigma_0/2=1.64$~fm$^2$, so the factor is about $3.6$. The ratio $\mathcal{N}_A/\mathcal{N}_p$ in the
dipole files at $\xpo=0.01$ and $\mathbf{b}=0$ is $3.6$ for $r_\perp=0.05$--$0.3$~GeV$^{-1}$, and
decreases for larger dipoles ($3.2$ at $1.3$~GeV$^{-1}$, $2.7$ at $2$~GeV$^{-1}$). The relation holds at
the initial condition only: after the BK evolution the ratio is smaller (Section 4).

The exclusive cross section is \emph{quadratic} in the amplitude, and for coherent production it is
integrated over the impact parameter:
\[
  \sigma_A\propto\int\mathrm{d}^2\mathbf{b}\,\mathcal{N}_A^2
  \approx A^2\left(\frac{\sigma_0}{2}\right)^2\mathcal{N}_p^2\int\mathrm{d}^2\mathbf{b}\,T_A^2\,,
  \qquad
  \sigma_p\propto\frac{\sigma_0}{2}\,\mathcal{N}_p^2\,.
\]
Therefore
\[
  R_{pA}\;\longrightarrow\;A\,\frac{\sigma_0}{2}\int\mathrm{d}^2\mathbf{b}\,T_A^2(\mathbf{b})\,.
\]
For a uniform sphere of radius $R_A$, $\int\mathrm{d}^2\mathbf{b}\,T_A^2=9/(8\pi R_A^2)$, so
\[
  R_{pA}\;\longrightarrow\;\frac{9}{8}\,\frac{A\,\sigma_0/2}{\pi R_A^2}\;\propto\;\frac{A}{A^{2/3}}=A^{1/3}.
\]
This is the well-known $A^{4/3}$ scaling of coherent production (the same as for coherent $J/\psi$):
$A^2$ from the amplitude squared, divided by the area $\pi R_A^2\propto A^{2/3}$ of the nucleus.

\textbf{Number for Pb.} With $A=208$, $\sigma_0/2=16.36$~mb and $R_A\approx7.1$~fm,
\[
  R_{pA}\approx\frac{9}{8}\,\frac{208\times1.64~\text{fm}^2}{\pi\,(7.1~\text{fm})^2}\approx2.4\,.
\]
This rough estimate agrees with our exclusive $R_{pA}$, which rises to about $2.3$ at
$p_\perp=12$~GeV. The full factor $A^{1/3}\approx5.9$ is never reached, because the nucleons are spread
over the large area of the nucleus. At low $p_\perp$ the dipoles are large, the nucleus is
saturated ($\mathcal{N}_A\to1$ cannot grow with $A$), and $R_{pA}$ drops towards $1$. Saturation is
only half of the reason: the $p_\perp$ and rapidity dependence is explained in Section 4.

\section*{2. Diffractive (TMD): $R_{pA}\to$ constant}

In the diffractive process the gluon momentum $K_\perp$ is integrated, up to $p_\perp$:
\[
  \frac{\mathrm{d}\sigma}{\mathrm{d}y\,\mathrm{d}^2\mathbf{p}\,\mathrm{d}\ln\xpo}
  =\mathcal{H}(z,p_\perp)\int^{p_\perp}\mathrm{d}^2\mathbf{K}\;\mathcal{G}_{\mathbb{P}}(x,K_\perp,\xpo)\,.
\]
The Pomeron gluon TMD $\mathcal{G}_{\mathbb{P}}$ is flat (saturated) for $K_\perp\lesssim Q_s$ and falls
like $Q_s^4/K_\perp^4$ above it (Ref.~[3], Sec.~4.3: plateau at low $K_\perp$, dilute tail
$\propto1/K_\perp^4$; unlike the quark DTMD, the gluon one has no Cronin peak). With the measure
$\mathrm{d}^2K=2\pi K_\perp\mathrm{d}K_\perp$, both the plateau and the tail give contributions of order
$Q_s^2$. So, even when $p_\perp$ is large, the $K_\perp$ integral is dominated by
$K_\perp\sim Q_s$: the saturation region, not small dipoles. The hard scale $p_\perp$ only enters
through $\mathcal{H}$, which is the same for both targets and cancels in $R_{pA}$.

The integrated distribution is then of order (area) $\times\,Q_s^2$:
\[
  \int\mathrm{d}^2\mathbf{K}\,\mathcal{G}_{\mathbb{P}}\;\propto\;\int\mathrm{d}^2\mathbf{b}\,Q_s^2(\mathbf{b})\,,
  \qquad Q_s^2(\mathbf{b})\propto A\,T_A(\mathbf{b})\,\frac{\sigma_0}{2}\,Q_{s,p}^2
  \;\Longrightarrow\;
  \int\mathrm{d}^2\mathbf{b}\,Q_s^2\propto A\,\frac{\sigma_0}{2}\,Q_{s,p}^2\,,
\]
which is $A$ times the proton value. Hence $R_{pA}\to$ constant, close to $1$.

In our figures it levels off at about $0.9$ rather than exactly $1$: the scaling argument holds only up
to the detailed shape of $\mathcal{G}_{\mathbb{P}}$ around $Q_s$, and our Pb dipoles (Glauber-sampled MVe
initial condition at each $b$, evolved with BK) are not exactly the proton one with
$Q_s^2\propto A\,T_A$. At low $p_\perp$ the upper limit $K_\perp\leq p_\perp$ cuts the integral below
$Q_{s,A}$, so Pb loses more than the proton and $R_{pA}$ drops to about $0.4$.

\section*{3. The sum}

The summed $R_{pA}$ is a weighted average of the two (an identity, not how it is computed),
\[
  R_{pA}^{\rm sum}=w\,R_{pA}^{\rm excl}+(1-w)\,R_{pA}^{\rm diff},\qquad
  w=\frac{\sigma_p^{\rm excl}}{\sigma_p^{\rm excl}+\sigma_p^{\rm diff}}\,.
\]
At very low $p_\perp$ the exclusive part dominates ($R\approx1$); from $p_\perp\approx1$~GeV the
diffractive TMD dominates (it falls like $p_\perp^{-4}$, the exclusive like $p_\perp^{-8}$), and the sum
follows its $R_{pA}$: a dip to about $0.7$, then a plateau at about $0.9$. The large exclusive
$R_{pA}$ at high $p_\perp$ does not show up in the sum, because there its weight is below $3\%$.

\section*{4. Exclusive: why $R_{pA}$ starts near 1, grows slowly and depends on $y$}

The numbers here are a recomputation of the exclusive $R_{pA}$ at fixed $q^+=2p^+$ directly from the
dipole tables (\texttt{data/Pb/mve/}, \texttt{data/proton/mve.dat}). It agrees with
\texttt{RpA\_fixed\_qp.py} within about $1\%$.

\textbf{The limit.} With the Woods--Saxon profile of the plotting script ($R=6.49$~fm, $a=0.542$~fm),
\[
  A\,\frac{\sigma_0}{2}\int\mathrm{d}^2\mathbf{b}\,T_A^2(\mathbf{b})=2.49\,,
\]
of which $0.05$ is the dilute edge beyond $b=32$~GeV$^{-1}$.

\textbf{Two effects keep $R_{pA}$ below the limit.} To separate them, the same calculation with the
dipole kept at its initial condition ($\xpo=0.01$, no evolution):
\begin{center}
\begin{tabular}{lcccccc}
\toprule
$p_\perp$ [GeV] & 0.5 & 1 & 2 & 4 & 8 & 12 \\
\midrule
no evolution       & 1.74 & 2.07 & 2.49 & 2.64 & 2.55 & 2.51 \\
with evolution, $y=0$ & 1.13 & 1.29 & 1.58 & 1.93 & 2.21 & 2.35 \\
with evolution, $y=2$ & 0.93 & 1.04 & 1.25 & 1.53 & 1.83 & 1.99 \\
\bottomrule
\end{tabular}
\end{center}
\begin{enumerate}
\item \emph{Saturation of large dipoles} (first row). Small $p_\perp$ means large dipoles, for which
$\mathcal{N}_A$ is no longer $A\,T_A\,\sigma_0/2$ times $\mathcal{N}_p$. This matters only below about
$2$~GeV: without evolution the limit $2.49$ is reached there.
\item \emph{BK evolution} (the other rows). The dipole is taken at
\[
  \xpo=\frac{2\,m_t\,e^{-y}}{\sqrt{s_{NN}}}\,,\qquad Y=\ln\frac{0.01}{\xpo}\,,
\]
so the evolution is longer for smaller $p_\perp$ and larger $y$: at small $p_\perp$, $Y\approx3$ for
$y=0$ and $Y\approx5$ for $y=2$; at $p_\perp=12$~GeV and $y=0$, $Y\approx0.8$. The evolution
lowers the ratio $\mathcal{N}_A/\mathcal{N}_p$ (see below). This is why $R_{pA}$ starts between $0.9$ and
$1.1$ (it is not a special limit), why it grows slowly up to $12$~GeV, and why it is smaller for larger $y$.
\end{enumerate}
The limit $2.49$ is reached only when $\xpo\geq0.01$, that is, for $m_t\gtrsim27\,e^{y}$~GeV, outside
the range of the figure.

\textbf{Why the evolution lowers the ratio.} The amplitude obeys the BK equation,
\[
  \frac{\partial\mathcal{N}(r,Y)}{\partial Y}=\int\mathrm{d}^2\mathbf{r}_1\,K(r,r_1,r_2)\,
  \big[\mathcal{N}(r_1)+\mathcal{N}(r_2)-\mathcal{N}(r)-\mathcal{N}(r_1)\mathcal{N}(r_2)\big],
\]
with the running-coupling kernel $K$. At $Y=0$, for small dipoles,
\[
  \mathcal{N}_A(r,\mathbf{b})=c(\mathbf{b})\,\mathcal{N}_p(r),\qquad
  c(\mathbf{b})=A\,T_A(\mathbf{b})\,\frac{\sigma_0}{2}=\frac{Q_{s,A}^2(\mathbf{b})}{Q_{s,p}^2}\,,
\]
with $c\simeq3.6$ at the center of Pb. Two things change with the evolution:
\begin{enumerate}
\item \emph{Anomalous dimension.} The amplitude of a small dipole is no longer proportional to
$r^2Q_s^2$ but to a smaller power,
\[
  \mathcal{N}(r,Y)\propto\big[r^2Q_s^2(Y)\big]^{\gamma},\quad\gamma<1
  \qquad\Longrightarrow\qquad
  \frac{\mathcal{N}_A}{\mathcal{N}_p}=\left[\frac{Q_{s,A}^2(Y)}{Q_{s,p}^2(Y)}\right]^{\gamma}<c\,.
\]
\item \emph{Running coupling.} The saturation scale grows as
$\mathrm{d}\ln Q_s^2/\mathrm{d}Y\propto\alpha_s(Q_s^2)$, which is smaller for the larger nuclear scale,
so the ratio $Q_{s,A}^2/Q_{s,p}^2$ itself decreases with $Y$.
\end{enumerate}
Since the exclusive cross section is quadratic in $\mathcal{N}$,
\[
  R_{pA}\simeq\frac{1}{A\,\sigma_0/2}\int\mathrm{d}^2\mathbf{b}
  \left[\frac{Q_{s,A}^2(\mathbf{b},Y)}{Q_{s,p}^2(Y)}\right]^{2\gamma},
\]
which is the limit above for $Y=0$, $\gamma=1$, and decreases with $Y$.

From the tables, for $r=0.2$~GeV$^{-1}$ at the center of Pb:
\begin{center}
\begin{tabular}{lccc}
\toprule
$Y$ & 0 & 2 & 4 \\
\midrule
$\gamma$ (slope of $\ln\mathcal{N}_A$ against $\ln r^2$) & 0.95 & 0.87 & 0.80 \\
$\mathcal{N}_A/\mathcal{N}_p$ & 3.6 & 3.0 & 2.6 \\
$3.6^{\gamma}$ (anomalous dimension alone) & 3.6 & 3.1 & 2.9 \\
\bottomrule
\end{tabular}
\end{center}
So most of the drop is the anomalous dimension. The rest at $Y=4$ ($2.9$ against $2.6$) corresponds to
$Q_{s,A}^2/Q_{s,p}^2$ falling from $3.6$ to about $3.3$. This split is an estimate at one dipole size
and one impact parameter.

\textbf{Remark on the diffractive TMD at low $p_\perp$.} The value $0.4$ of Section 2 is a ratio of
areas. For $K_\perp\leq p_\perp\ll Q_s$ the gluon dipole is black in both targets, the cross section is
proportional to the transverse area, and
\[
  R_{pA}\to\frac{S_A}{A\,\sigma_0/2}\approx0.42
  \quad\text{for}\quad S_A\approx143~\text{fm}^2=\pi\,(6.7~\text{fm})^2 ,
\]
independently of rapidity, as in the figure. In the exclusive process this does not happen, because
the size of the $q\bar q$ dipole is limited by $1/m_c$ even at $p_\perp\to0$. (This is the black-disk
limit of the formula; it was not recomputed numerically.)

\section*{5. Summary}

\begin{center}
\begin{tabular}{llll}
\toprule
Process & dominated by & $\sigma_A$ & $R_{pA}$ at large $p_\perp$ \\
\midrule
Exclusive & small dipoles, dilute, quadratic in $\mathcal{N}$ & $\propto A^{4/3}$ & $\propto A^{1/3}$, $\approx2.4$ for Pb \\
Diffractive (TMD) & $K_\perp\sim Q_s$, saturation & $\propto A$ & $\approx0.9$--$1$ \\
\bottomrule
\end{tabular}
\end{center}

A possible sentence for the paper: \emph{The exclusive contribution at large $p_\perp$ probes small
dipoles in the dilute regime and scales as $A^{4/3}$, giving $R_{pA}\propto A^{1/3}$; the diffractive
TMD contribution, instead, is dominated by gluon transverse momenta of the order of the saturation scale
even at large $p_\perp$, so that it scales approximately as $A$ and remains sensitive to saturation.}

\end{document}
